\chapter{Desarrollo}\label{cap:Desarrollo}
% Índice v4, cap. 4: aplica el ciclo del cap. 3 al Cuaderno de Campo. Cada
% apartado muestra artefactos, iteraciones con la IA y decisiones del alumno.
% Se redacta a partir de docs/proceso/registro-iteraciones.md, no de memoria.

\section{Constitución y arquitectura}\label{sec:Constitucion}
% constitution.md (v1.0.0, 2026-10-23) y la justificación de cada principio;
% arquitectura (frontend PWA, API FastAPI, PostgreSQL + pgvector, RAG, servicios
% externos) con diagrama de componentes; despliegue local con Docker + demo.
\pendiente{constitución, arquitectura y despliegue}

\section{Fase de ideación con Claude: especificar y clarificar}\label{sec:Ideacion}
% Iteraciones hasta spec.md; REQ-xx, historias y criterios de aceptación,
% incluidos los del RAG. El MD completo va en el anexo A.
\pendiente{especificación y clarificación}

\section{Fase de generación con Claude Code en VS Code}\label{sec:Generacion}
% Preparación (CLAUDE.md, Spec Kit, hooks); plan.md y tasks.md con correcciones;
% implementación en bucle; estructura resultante y construcción del RAG.
\pendiente{planificación, tareas e implementación}

\section{Fase de despliegue y testing}\label{sec:Testing}
% Pipeline CI/CD; E2E con Playwright y tabla REQ -> prueba; rendimiento con k6;
% evaluación del RAG (fidelidad y relevancia).
\pendiente{pipeline, pruebas y evaluación del RAG}

\section{Resultados del proceso}\label{sec:ResultadosProceso}
% Métricas del registro: iteraciones, código corregido, tiempos, cobertura y
% resultados de pruebas, con su análisis.
\pendiente{métricas del proceso}
